\chapter{Zero-Dimensional Cellular Automata}

\section{Preliminaries}

A two-dimensional lattice may be thought of as the points of the plane
with integer coordinates, and a one-dimensional lattice as the integer
points of the real line.  A zero-dimensional lattice, then, consists of
a single point --- not much of a lattice at all.  The parallelism
inherent in higher-dimensional CA is totally lacking, so one might well
suppose the 0-dimensional case trivial and devoid of interest.  But no
one would call the concept of \emph{function} trivial, and discrete
functions of one variable, regarded as dynamical systems, \emph{are}
0-dimensional CA.

We now take a little detour through the land of finite dynamical
systems, illustrating important concepts these simple systems share
with far more complex ones, from higher-dimensional CA to systems of
differential equations.  To fit the CA framework we must specify a
lattice, the allowable cell states, neighborhoods, transition rules,
and an initial state.  The lattice is a single cell; the allowable
states are the integers modulo some fixed integer $n$; the neighborhood
of our lonely cell is the cell itself; the transition rule is an
arithmetic function $f$ whose output we reduce modulo $n$; and the
initial state is a single integer $X_0$.  There are only $n$ possible
initial states: $0,1,\dots,n-1$.  Letting the dynamics unfold, we
iterate the recursion
\begin{equation}
  X_{k+1} = f(X_k) \pmod n,
\end{equation}
obtaining the sequence
\[
  X_0,\, X_1,\, X_2,\, \dots,\, X_M,\, \dots,\, X_N,\, \dots\;.
\]

This sequence is termed a \emph{trajectory} --- the same word used for
the path traversed by an initial point under the influence of a
differential equation.  Here the trajectory is a discrete sequence
rather than a continuous path, and a recursive arithmetic function
replaces the differential equation (an ``infinitesimal recursion
rule'').  The near-identical vocabulary reflects the closely analogous
relationship of the two realms.

Since each $X_i$ can take only one of $n$ values, at some point ---
say at $X_M$ --- we must encounter a value for the second time; from
then on we repeat the sequence that followed its first occurrence.  We
have entered a \emph{cycle}.  The portion of a trajectory preceding its
entry into a cycle is a \emph{transient}, and a cycle possessing
transients is an \emph{attractive cycle}.  The analogous concept for
differential equations is the \emph{limit cycle}.

Trajectories give a local way of understanding dynamics.  For a global,
qualitative understanding of differential equations,
Poincar\'e introduced the \emph{state space} (or \emph{phase space});
the analogous concept in the discrete realm is the \emph{basin of
attraction field}.  These are easy to construct in the 0-dimensional
case and they illustrate the properties of basin-of-attraction fields
in general.

\section{Example: $x^2+1 \pmod{17}$}
\label{sec:mod17}

For our first example let $n=17$ and define the recursion function by
\begin{equation}
  f(x) = x^2 + 1 \pmod{17}.
\end{equation}
Starting with $x=0$ we get the trajectory
$0, 1, 2, 5, 9, 14, 10, 16, 2, \dots$.  Here $2$ has appeared a second
time, and so begins the attractive cycle $2, 5, 9, 14, 10, 16$.  The
module \texttt{cadyn.fds} computes the basin-of-attraction field shown
in Figure~\ref{fig:mod17}.

\begin{figure}[htbp]\centering
\includegraphics[width=0.8\textwidth]{ch2_basin_mod17}
\caption{Basin of attraction field for $f(x)=x^2+1 \pmod{17}$.  Cycle
edges are red, transient edges green.}
\label{fig:mod17}
\end{figure}

Notice the attractive cycle and the incoming transient branches.  For
this function there is only one attractive cycle; ordinarily there are
several, and the basin-of-attraction field consists of the entire
collection of cycles with their associated transients.  Each node has
exactly one outgoing edge: the system is \emph{deterministic}.  Some
nodes have more than one incoming edge: the system is
\emph{irreversible}, and the entropy of the system tends to decrease as
it evolves.  CA of all dimensions share these qualities, and the
tendency of entropy to decrease in CA is no doubt partly responsible
for the emergent behavior so often evident in such systems.

Notice too that the two predecessors of any node with two predecessors
always sum to $17$.  This follows from the algebraic identity
$x^2 = (-x)^2$ together with the modular-arithmetic fact
$-x \equiv n-x \pmod n$: hence
\[
  x^2 \equiv (n-x)^2, \qquad x^2+1 \equiv (n-x)^2+1,
\]
and so $f(x) \equiv f(n-x) \pmod n$ for $0 \le x \le n-1$.  And indeed
$x + (n-x) = n$.  Here $n=17$, but the same holds for the later
examples using this function; the reader should verify it in
Sections~\ref{sec:mod37} and~\ref{sec:mod221}.

The outermost states in Figure~\ref{fig:mod17}, with values
$4, 6, 7, 8, 11, 12, 13, 15$, have no predecessor states: after one
generation the CA is guaranteed not to occupy any of them.  Such
outlying states are known as \emph{Garden of Eden} states.

Before further examples we digress to define the important concept of
entropy.  Entropy is not separate from information, in precisely the
way the concept of a half-empty glass is not separable from that of a
half-full one; out of this concept grows the entire field of
information theory.

\section{Shannon's Entropy}
\label{sec:entropy}

Entropy for a 0-dimensional CA with $n$ states may be defined via
Shannon's formula
\begin{equation}
  H_t \;=\; -\sum_{k=0}^{n-1} p_t(k)\,\log_2 p_t(k),
\end{equation}
where $p_t(k)$ is the probability that the system has value $k$ after
$t$ time steps from a uniformly random initial value, and where we
agree that $0\log 0 = 0$.  Using base-2 logarithms measures entropy in
\emph{bits}.  $H$ measures the system's basic information potential.

Two extreme cases illuminate the formula.  First, consider a CA in
which every initial state evolves to $0$ (e.g.\ $f(x)=0$).  For $t$
large enough, $p_t(0)=1$ and all other probabilities vanish, so the sum
reduces to $\log_2 1 = 0$: zero entropy.  Entropy connotes randomness
and uncertainty, and here we are \emph{a priori} certain of the
outcome.  Second, consider a CA in which all states remain equally
likely for all time (e.g.\ $f(x)=x+1$).  Then $p_t(k)=1/n$ for all $k$,
and the formula gives
\begin{equation}
  H = -\log_2 \tfrac1n = \log_2 n,
\end{equation}
the maximum entropy possible for a 0-dimensional CA with $n$ states
(Table~\ref{tab:maxent}).  Note that $H$, rounded up, is the number of
binary digits needed to represent $n$.

\begin{table}[htbp]\centering
\begin{tabular}{l|ccccccccc}
\toprule
$n$ & 1 & 2 & 5 & 10 & 15 & 20 & 30 & 50 & 100\\
\midrule
$S=\log_2 n$ & 0 & 1 & 2.32 & 3.32 & 3.91 & 4.32 & 4.91 & 5.64 & 6.64\\
\bottomrule
\end{tabular}
\caption{Maximum entropy $S$ as a function of the number of cell
states $n$.}
\label{tab:maxent}
\end{table}

To compute $p_t(k)$, the software counts the number of initial values
that end up with value $k$ after $t$ iterations; call this
$\mathrm{Num}(k,t)$.  Since there are $n$ equally likely initial
states,
\begin{equation}
  p_t(k) = \frac{\mathrm{Num}(k,t)}{n}.
\end{equation}
$H_0$ is thus always the maximal entropy of the system, and once every
value has settled into a cycle the entropy remains constant at the
system's minimum.\footnote{Strictly, once all transients have died out
the multiset of occupied states is permuted by $f$, so the entropy is
exactly constant from then on.}  In general, the smaller the entropy,
the greater one's \emph{a priori} information about the state of the
system: entropy measures one's uncertainty about a system's state.

\section{Example: $x^2+1 \pmod{37}$}
\label{sec:mod37}

Figure~\ref{fig:mod37} depicts the basin-of-attraction field for
\begin{equation}
  f(x) = x^2 + 1 \pmod{37}.
\end{equation}
There are three separate attractive cycles: two of length~1 (the fixed
points $27$ and $11$) and one of length~2 (the cycle $8, 28$).
Attractive 1-cycles are usually just called \emph{attractors} or
\emph{fixed points}.

\begin{figure}[htbp]\centering
\includegraphics[width=0.9\textwidth]{ch2_basin_mod37}
\caption{Basin of attraction field for $f(x)=x^2+1\pmod{37}$: the
attractor 27 with a basin of 8 states, the 2-cycle $\{8,28\}$ with a
basin of 6 states, and the attractor 11 with a basin of 23 states.}
\label{fig:mod37}
\end{figure}

Since the system eventually evolves to one of four values, the entropy
should be relatively small after six iterations.
Figure~\ref{fig:entropy37} graphs the entropy $H$ as a function of the
number of iterations.

\begin{figure}[htbp]\centering
\includegraphics[width=0.65\textwidth]{ch2_entropy_mod37}
\caption{Entropy as a function of time for $f(x)=x^2+1 \pmod{37}$.}
\label{fig:entropy37}
\end{figure}

The system ends up at one of its attractors ($27$ or $11$) or on its
2-cycle $(8,28)$ after 6 time steps, so thereafter it occupies one of
these 4 values.  A uniformly random system with 4 states has entropy
$\log_2 4 = 2$, rather larger than this system's minimum entropy,
$1.478$.  Our smaller final entropy is due mostly to the large number
of transients feeding attractor 11: with 23 states in its basin, the
probability of ending at 11 is $23/37 = 0.621\ldots$, so most of the
time the system ends there.  This additional information means less
uncertainty than in the uniform case, where the probability of ending
at 11 would be only $1/4 = 0.25$; hence the lower entropy.  The greater
one's \emph{a priori} information about a system, the less one's
uncertainty about its states, and the lower its entropy.

\section{Example: $x^2+1 \pmod{221}$}
\label{sec:mod221}

Figure~\ref{fig:mod221} depicts the basin-of-attraction field for
\begin{equation}
  f(x) = x^2 + 1 \pmod{221}.
\end{equation}
In this case there are four separate attractive cycles: two of
length~12 and two of length~6, accounting for the 36 cyclic states.

\begin{figure}[htbp]\centering
\includegraphics[width=\textwidth]{ch2_basin_mod221}
\caption{Basin of attraction field for $f(x)=x^2+1\pmod{221}$: two
attractive 12-cycles and two attractive 6-cycles, with their
transients.  Note the topological repetition between basins of equal
cycle length.}
\label{fig:mod221}
\end{figure}

The topological repetition observed among basins of attraction is very
common for large moduli, and the redundancy becomes even more
pronounced for CA on higher-dimensional lattices.\footnote{For this
example the repetition has a tidy explanation: $221 = 13\cdot17$, and
by the Chinese Remainder Theorem the system is the direct product of
$x^2+1$ mod 13 and mod 17, whose cycle structures combine to produce
the paired 6- and 12-cycles.}
Figure~\ref{fig:entropy221} shows the entropy graph for this system.
The outcome of many iterations is far less certain than for the system
of Section~\ref{sec:mod37}, so the minimum entropy is higher
($H_{\min}=4.983$ versus $1.478$).

\begin{figure}[htbp]\centering
\includegraphics[width=0.65\textwidth]{ch2_entropy_mod221}
\caption{Entropy as a function of time for $f(x)=x^2+1 \pmod{221}$.}
\label{fig:entropy221}
\end{figure}

The system reaches one of its 4 cycles after 3 time steps and is
thereafter confined to the 36 cyclic values.  A uniformly random system
with 36 states has entropy $\log_2 36 = 5.170$, very close to this
system's minimum entropy.  Our slightly smaller final entropy is due to
the greater number of transients associated with the 12-cycles, which
makes ending in a 12-cycle more probable than ending in a 6-cycle; this
extra information lowers the entropy below that of a uniformly random
system.

\section{Example: $x^3 + 2x^2 + 3x + 2 \pmod{28}$}
\label{sec:cubic}

Figure~\ref{fig:cubic} shows the basin-of-attraction field for
\begin{equation}
  f(x) = x^3 + 2x^2 + 3x + 2 \pmod{28},
\end{equation}
whose basin-of-attraction field comprises two basins: an attractive
6-cycle and an attractive 2-cycle, each with associated transients.
Figure~\ref{fig:entropycubic} shows the entropy graph.

\begin{figure}[htbp]\centering
\includegraphics[width=0.8\textwidth]{ch2_basin_cubic_mod28}
\caption{Basin of attraction field for
$f(x)=x^3+2x^2+3x+2 \pmod{28}$: an attractive 6-cycle and an
attractive 2-cycle.}
\label{fig:cubic}
\end{figure}

\begin{figure}[htbp]\centering
\includegraphics[width=0.65\textwidth]{ch2_entropy_cubic_mod28}
\caption{Entropy as a function of time for
$f(x)=x^3+2x^2+3x+2 \pmod{28}$.}
\label{fig:entropycubic}
\end{figure}

\section{Ulam's Function}
\label{sec:ulam}

In the late 1940s Stanis\l{}aw Ulam introduced the $3N+1$ problem to
the Los Alamos scientific community (though he didn't invent it), which
was beginning to play with the very first digital computers.  The
problem involves a simple recursively defined arithmetic function that
to this day holds unfathomed mysteries:
\begin{equation}
  f(n) =
  \begin{cases}
    n/2, & \text{if } n \text{ even};\\
    3n+1, & \text{if } n \text{ odd}.
  \end{cases}
\end{equation}
The observation Ulam circulated is that no matter what positive integer
one starts with, iterating $f$ appears always to arrive at the
attractive cycle $4\to2\to1$.  Nobody knows whether \emph{all} positive
integers belong to this cycle's basin of attraction.

\expand{The conjecture (also attached to the names Collatz, Kakutani,
Syracuse, and others) remains open.  The first edition noted
verification for all $n < 10^{12}$; as of the mid-2020s it has been
verified for all $n < 2^{71} \approx 2.36\times10^{21}$ (Barina 2020
and subsequent distributed computation), and Terence Tao proved in 2019
that ``almost all'' orbits attain almost bounded values --- yet a full
proof remains out of reach, vindicating Erd\H{o}s's judgment that
``mathematics is not yet ripe for such problems.''}

If zero is allowed as an argument, it becomes another attractor.  The
module \texttt{cadyn.fds} has this function built in, modulo arbitrary
values, and when computed with modular arithmetic zero has its own
basin of attraction.  Figure~\ref{fig:ulam} shows the computation
modulo 20 and modulo 40, and Figure~\ref{fig:ulam47} the computation
modulo 47.  Notice the surprising occurrence
of yet another attractor, $23$, in the modulus-47 case.

\begin{figure}[htbp]\centering
\includegraphics[width=0.88\textwidth]{ch2_basin_ulam_mod20}\\[1.5ex]
\includegraphics[width=0.88\textwidth]{ch2_basin_ulam_mod40}
\caption{Ulam's function modulo 20 (top) and 40 (bottom).}
\label{fig:ulam}
\end{figure}

\begin{figure}[htbp]\centering
\includegraphics[width=0.95\textwidth]{ch2_basin_ulam_mod47}
\caption{Ulam's function modulo 47, with the surprise attractor
at 23.}
\label{fig:ulam47}
\end{figure}

\begin{figure}[htbp]\centering
\includegraphics[width=0.62\textwidth]{ch2_entropy_ulam_mod47}
\caption{Entropy history for Ulam's function modulo 47.}
\label{fig:ulam-entropy}
\end{figure}
